\section{问题重述}
长江流域淡水渔业发达，生物多样性丰富，长期过度捕捞、工业废水、水域塑料污染以及葛洲坝、三峡大坝等水利工程也在这里叠加。洄游通道受阻后，渔业资源持续衰退，典型食物链随之退化。长江流域自2021年1月1日实施十年禁渔，局部水域水质和水生植被已有恢复；2025年监测数据显示，青、草、鲢、鳙四大家鱼资源量约为禁渔前的1.8倍，表明禁渔已带来资源恢复。

生态系统的恢复并非总是朝着良性方向发展。部分水域鱼类数量增长过快，对水生植物和浮游动物的过度摄食导致水体浑浊、底栖植被减少，甚至出现"水下荒漠"。长江江豚数量已回升至约1249头，2024年又发现长江鲟自然繁殖产卵场；鳄雀鳝、巴西龟等外来物种则已在鄱阳湖等支流湖泊定殖。十年禁渔后的系统未必更加健康，营养级失衡、连通性不足与入侵扩散同时存在，因此有必要分别衡量"资源恢复"和"系统健康"。\cite{report2022,report2023}。

题目提出的五个子问题依次为：
\begin{itemize}
\item 问题一：分析禁渔政策对水生植物、浮游动物及四大家鱼种群动态的影响。
\item 问题二：研究禁渔产生的生态收益能否向上传递至长江鲟和长江江豚等顶级捕食者。
\item 问题三：构建长江流域简化鱼类食物链模型并探讨其长期演化规律。
\item 问题四：判别三级食物链系统在不同承载力条件下的动力学行为（稳定、周期振荡或混沌）。
\item 问题五：将工业废水排放、塑料污染和外来物种入侵等胁迫因子纳入统一评价框架，综合评估十年禁渔对长江渔业生态系统的整体效应，并提出针对性的治理对策建议
\end{itemize}
全文的时间口径统一为：2021年作为政策启动基准年，2022年采用公报基准，2025年用于对照恢复效果\cite{yangtzeA}。

